\section{Forma del modelo, eficiencia de inferencia y campo receptivo efectivo}

Una vez alcanzado cierto nivel de accuracy, el siguiente paso no consiste en ampliar a ciegas todas las dimensiones. La depth, la width, la MLP width, el número de heads y el patch size del Transformer modifican la optimización, la activation memory, la eficiencia de paralelización y la latencia de inferencia. La clase lo explica con tres análisis: shape, speed y receptive field. Con los mismos parámetros o el mismo compute, la elección de la forma puede cambiar los resultados de manera considerable.

\subsection{El uniform scaling no siempre es la mejor opción}

La \term{model shape} es la manera en que los parámetros se distribuyen entre la profundidad, la hidden width, la MLP width, las heads y la patch resolution. Ampliar todas las dimensiones por igual es sencillo, pero puede destinar el presupuesto a direcciones con un rendimiento marginal bajo. En la figura, cada color representa una forma distinta de scaling y se compara su desempeño promedio en función del relative compute.

\lecturefigure{slide-23-transformer-shape.jpg}{La depth, la width, el patch size y la MLP width de ViT influyen de manera distinta en la compute efficiency.}{Video oficial, 46:20.}

\begin{knowledgebox}{Lectura de la figura: que las curvas estén próximas no significa que los sistemas sean equivalentes}
El eje vertical es la puntuación promedio y el horizontal, el relative compute. Una curva más alta indica mayor exactitud con el mismo presupuesto. El depth scaling puede aumentar los niveles de representación, pero alarga el camino secuencial; la width favorece la paralelización, pero amplía la activation; un patch más pequeño aumenta la resolución espacial, pero dispara el costo de la attention. La shape final debe considerar a la vez el throughput de entrenamiento, el batch de inferencia, la memory y la resolución objetivo.
\end{knowledgebox}

\subsection{Inference speed: las constantes de complejidad y la correspondencia con el hardware pueden invertir la clasificación}

Los FLOPs teóricos solo describen la cantidad de operaciones; la inference speed real depende también de la kernel implementation, el batch size, la memory bandwidth y la accelerator shape. Las dense matrix operations de ViT aprovechan con facilidad el hardware moderno, pero el aumento del número de tokens amplía con rapidez la attention y la activation.

\lecturefigure{slide-24-inference-speed.jpg}{Variación del throughput de inferencia de distintas variantes de ResNet y ViT según el tamaño de la imagen o del modelo.}{Video oficial, 47:56.}

\begin{knowledgebox}{Lectura de la figura: las curvas de velocidad requieren hardware context}
El eje vertical representa la velocidad o el throughput de inferencia y el horizontal corresponde al tamaño de la entrada o a la configuración del modelo. Al comparar, deben fijarse el accelerator, la precision, el batch size, el compiler y el preprocessing. Un modelo que lidera en throughput con batch grande puede quedar atrás en latencia con batch-1; que sea favorable para el entrenamiento en investigación tampoco implica que lo sea para dispositivos móviles.
\end{knowledgebox}

\begin{warningbox}{El patch size es un parámetro oculto del sistema}
Reducir $P$ conserva información de grano fino, pero hace crecer $N=HW/P^2$. Esto afecta al mismo tiempo los FLOPs de attention, la memoria de KV y de activation, la tensor shape después de la carga de datos y la resolución de la dense prediction posterior. Reportar «ViT-B» sin indicar `/16` o `/32` omite una variable de costo esencial.
\end{warningbox}

\subsection{Effective receptive field: la attention también forma una estructura jerárquica}

En teoría, cada capa de ViT permite interacción global, pero en la práctica las attention heads pueden ser más locales en las capas superficiales y ampliar su alcance en las profundas. El \term{effective receptive field} es la dependencia real del modelo respecto de distintas distancias espaciales, no el alcance teóricamente posible. Ayuda a determinar si la estructura local está codificada de forma fija o surge de los datos.

\lecturefigure{slide-25-receptive-field.jpg}{La distancia media de atención de distintas capas y attention heads muestra la estructura de lo local a lo global que aprende ViT.}{Video oficial, 49:20.}

\begin{knowledgebox}{Lectura de la figura: la nube de puntos muestra la heterogeneidad entre heads}
El eje horizontal es la layer depth y el vertical, la attention distance promedio; cada punto puede corresponder a una head. En las capas superficiales hay heads locales, y en las profundas la mayoría de las heads abarca distancias mayores, lo que indica que el modelo forma una jerarquía. Sin embargo, el attention weight no equivale a la feature attribution ni basta por sí solo para demostrar la función causal de una head; se requieren experimentos de masking o de intervention para verificarlo.
\end{knowledgebox}

\teachervoice{En esta parte, el docente vuelve una y otra vez a la pregunta de «cómo se define una misma escala». El número de parámetros, los FLOPs, el throughput real y la capacidad de representación no coinciden. En la práctica, conviene desglosar el nombre del modelo en depth, width, patch, resolution y precision, para no ocultar las diferencias entre sistemas detrás de un solo número total de parámetros.}

\subsection{Resumen de la sección}

Escalar un modelo es un problema de asignación de recursos. La shape determina cómo se distribuye el cómputo entre depth, width y token resolution; el hardware determina cómo los FLOPs se traducen en velocidad, y el campo receptivo efectivo permite revisar la jerarquía espacial que el modelo aprende en realidad. Los tres factores determinan en conjunto si vale la pena un «ViT más grande».
